%% Use the "review" option when submitting for review.
%% Removing the "review" option will switch the
%% manuscript to a two-column layout.
%%
%% Use option "jog" for submissions to the Journal of
%% Glaciology; use "aog" for submission to the Annals
%% of Glaciology.
% \documentclass[review,aog]{igs}
%\documentclass[,aog]{igs}
\documentclass[review, jog]{igs}
\usepackage[utf8]{inputenc}
\usepackage{natbib}
\usepackage{mathabx}
\usepackage{graphicx}
\usepackage{siunitx}
\usepackage{url}
\usepackage[final]{changes}
% if you really must have numbered sections, remove
% the % from the beginning of the following command
% and insert the level of sections you wish to be
% numbered (up to 4):

% \setcounter{secnumdepth}{2}

\jourvolume{12}
\jourissue{1}
\jourpubyear{2026}

\begin{document}

\title[Floe-scale sea ice motion]{Characterizing the Marginal Ice Zone in the Greenland Sea Through Seasonal Floe-Scale Sea Ice Observations}

\author[Watkins and others]{Daniel M. WATKINS,$^1$ 
  Ellen M. BUCKLEY,$^2$ Minki KIM,$^1$ 
  Jennifer K. HUTCHINGS,$^3$
  Monica M. WILHELMUS$^1$}

\affiliation{%
$^1$Center for Fluid Mechanics, Brown University, Providence, RI, USA\\
$^2$Department of Earth Science \& Environmental Change, University of Illinois Urbana-Champagne, Urbana, IL, USA\\
$^3$College of Earth, Ocean, and Atmospheric Sciences, Oregon State University, Corvallis, OR, USA \\
  Correspondence: Monica M. Wilhelmus 
  \email{mmwilhelmus@brown.edu}}

\begin{frontmatter}
\maketitle

\begin{abstract}
Sea ice motion in the marginal ice zone (MIZ) is inherently multi-scale, influenced by the relationship between the floe size distribution (FSD) and the temporal and spatial scales of winds and ocean currents. Joint observations of the FSD and motion are necessary to develop and validate models of the MIZ. However, such observations are rare.
We present observations of floe-scale Lagrangian sea ice translation and rotation in the East Greenland MIZ spanning the spring-to-summer transition for years 2003-20.
The observations are derived from moderate-resolution optical imagery using the Ice Floe Tracker (IFT) algorithm and consist of floe shapes, geometric properties, \added{reflectance}, and trajectories. We investigate FSD seasonality, links between floe size, translation, and rotation, and length scale dependence of horizontal deformation. The FSD evolves from April to July, with increasing prevalence of small floes. Variability of both drift speed and rotation depend on the floe size. \added{Deviations from the mean drift resolved by IFT are exponentially distributed, suggesting ocean turbulence is resolved.} Deformation rates show strong length scale dependence and vary seasonally. \deleted{Comparisons to large-scale sea ice motion and to discrete element model simulations show that deviations from the mean drift resolved by IFT are exponentially distributed, suggesting ocean turbulence is resolved.} These observations provide new information for sea ice model development and validation in the MIZ.
\end{abstract}
\end{frontmatter}

\section{Introduction}
% Why the floe scale matters
The size of an individual plate of ice--the \emph{floe scale}--is foundational to sea ice dynamics in the marginal ice zone (MIZ).
The relationship between the floe scale and the spatial scales of atmospheric and oceanic motion is important for understanding coupled interactions, as stresses from the atmosphere and ocean are integrated across the surface of the floe \citep{brenner2023_ScaleDependentAirSea,
gupta2024_EddyInducedDispersion,
kim2025_CharacterisationSea}. The spacing between floes and the range of floe sizes in an area, characterized by the sea ice concentration (SIC) and the floe size distribution (FSD), respectively, strongly affect air-ocean fluxes, lateral melt rates, wave-ice interactions, and floe-floe interactions \citep{Mcnutt2003, loose2014_ParameterModel, roach2018_EmergentSea,
wenta2018_InfluenceSpatial, horvat2022_FloesMarginal, herman2022_GranularEffects, manucharyan2022_SpinningIce, brenner2024_ScalingSimulations}.
Thus, joint observations of floe shapes and motion are needed to characterize MIZ sea ice dynamics.

Spatial and temporal variability in ice motion is directly related to the FSD and SIC \citep{overland1995_HierarchySea, dumont2022_MarginalIce,
herman2022_GranularEffects}.
The FSD is constantly evolving as floes either become larger due to pieces of ice interlocking, rafting, ridging and fusing, or smaller  due to fracture and melt \citep{zhang2015_SeaIce, horvat2016_InteractionSea, bateson2020_ImpactSea, roach2024_PhysicsSeasonal}.
The tight coupling between the seasonally evolving FSD, mesoscale and submesoscale ocean currents, waves, and wind stress in the MIZ results in complex sea ice motion patterns with low spatial correlation (meters to kilometers) and short temporal scales (hours to days) \citep{thorndikeSeaIceMotion1982, hakkinen1987_FeedbackIce, johannessen1987_MesoscaleEddies,
johannessen1987_IceEdgeEddies, feltham2005_GranularFlow,  cole2017_IceOcean, watkins2023_EvidenceAbrupt}.

\added{The dynamical complexity of the MIZ}\deleted{This complexity} poses a key challenge for sea ice observation systems\deleted{and modeling.}
\deleted{Sea ice motion is routinely monitored through drifting buoys and by processing remote sensing imagery }\citep{webster2022_ObservingArctic, gerland2019_EssentialGaps, sandvenSeaIceRemote2023}.
The physical processes that can be resolved are constrained \deleted{both} by the spatial resolution and acquisition frequency of observations \added{as well as whether the system is designed in a Lagrangian or Eulerian framework. }
\deleted{as well as the reference frame of observations: buoys capture ice motion in a Lagrangian sense, whereas most remote sensing products produce Eulerian motion estimates.}
\deleted{Outside dedicated experiments where observations sources are merged (Uttal and others, 2002; Hwang and others, 2017; Rabe and others, 2022) few observation sources exist that jointly observe ice dynamics and floe properties.}
\added{On the one hand, }buoys track the motion of individual ice floes with high precision. \deleted{At large scales (greater than 100 km)}
\added{At scales greater than 100 km,} operational buoy deployments through the International Arctic Buoy Program (IABP, \cite{rigor2008}) provide long-term coverage of the sea ice cover. 
Due to the difficulty and expense of deployment, \added{however}, this data is spatially and temporally sparse, particularly in the MIZ \citep{gerland2019_EssentialGaps, brunette2022_NewStatedependent}.
% To add here: buoy data driving model development, including observing scale dependence
In addition, buoy-derived measurements are limited to point-wise \deleted{measurements}\added{acquisitions, which must be collated with other sensors to estimate sea ice properties.}\deleted{measurements. Hence the representativeness of a buoy observation in space must be estimated, for example, through collation with other sensors}.
\added{For example, the spatiotemporal scale dependence in sea ice motion and deformation is obtained using measurements from buoy arrays \citep{Hutchings2011, itkin2017_ThinIce}.}
\added{On the other hand,} remote sensing observations from satellites and aircraft enable observation of sea ice properties across large spatial regions.
Ice motion is computed from satellite imagery by \added{cross-correlating pixel intensity between consecutive image pairs}\deleted{estimating the displacement of sea ice from the shift in pixel intensity from one image to the next, for example through maximizing cross-correlation} \citep{ninnis1986_AutomatedExtraction, kwok1998_SeaIce}.
\deleted{For such methods, the motion vectors produced}
\added{The produced motion vectors} represent an area-average within the spatial window used by the correlation algorithm.
In the pack ice, where spatial correlation is high, the sea ice moves coherently and displacement estimates have high confidence; however, the same spatial window in the MIZ may include thousands of jostling floes, resulting in higher uncertainty \citep{sumata2014_IntercomparisonArctic, gui2020_ValidationRemotesensing}. 
Moreover, many products are only available seasonally, with few products providing estimates of motion during the summer melt season due to the presence of liquid water at the snow-and-ice surface affecting the quality of the images and the ability to distinguish between surface melt water and leads \citep{lavergne2010_SeaIce, girard-ardhuin2012_EnhancedArctic}. 
Thus the representativeness, availability, and uncertainty of remote sensing-derived ice motion vectors varies spatially and temporally, and is degraded in the MIZ.


In contrast to sea ice motion, the FSD is not routinely measured, though standard operational ice charts do include information on predominant floe size categories \citep{dedrickUSNationalNaval2001, afanasyevaAARIMethodologySea2019}.
\added{
Furthermore, to our knowledge, no existing data products provide the FSD and ice motion simultaneously.}
\deleted{To date, knowledge of the properties and variability of the FSD comes through individual studies.}
\deleted{Observation of the sea ice FSD from airborne remote sensing imagery dates back to Rothrock and Thorndike, 1984. Noting an approximately linear FSD probablity density function (PDF) in log space, they suggested characterization of the FSD via}
\added{
Traditional methods and image types used to characterize the FSD vary \citep{pagetDeterminingFloesizeDistribution2001,  toyota2006_CharacteristicsSea, hwang2013_IntercomparisonSatellite}. Most commonly, the FSD is characterized via} a power law with the form
\begin{linenomath*}
\begin{equation}
    p(x) = cx^{-\alpha} \label{powerlaw}
\end{equation}
\end{linenomath*}
where $x$ is a measure of floe size, $c$ is a normalization constant and the parameter $\alpha$ is the slope of the power law distribution on a log-log scale. This representation assumes an approximately linear FSD probability density function (PDF) in log-log space  \added{\citep{rothrock1984_MeasuringSea}}. 
\deleted{%moved
Numerous methods and image types have been used to characterize the FSD (Paget and others, 2001; Toyota and others, 2006; Hwang, 2013).}
\deleted{As a result, measured FSD properties (and methods) vary substantially across studies (Stern and others, 2018) often showing a clear departure from Eq. 1}
\deleted{Note that the form of} \added{While this form is convenient given its direct application in models, the power law in} Eq. (1) is an empirical, not theoretical result \citep{herman2010_SeaiceFloesize}, and distinct mechanisms governing the evolution of the FSD may lead to other forms of the FSD PDF. 
\added{Indeed, FSD characteristics vary substantially across studies \citep{stern2018_ReconcilingDisparate}, often showing a clear departure from Eq. \ref{powerlaw}.
Differences in methodology across FSD studies poses a challenge for developing a unified theory for FSD evolution.
}
For example, theoretical and model-based analyses of floe breakup through wave action suggest a log-normal FSD \citep{montiel2022_TheoreticalFrameworka, mokus2022_WavetriggeredBreakup}, \added{while \cite{toyotaSizeDistributionShape2011} found that observed FSDs had different power law slope parameters in different size regimes.}


\deleted{
Development and validation of models capable of simulating small-scale sea ice variability requires observations of sea ice across scales.
Traditional models of sea ice dynamics and thermodynamics rely on a continuum approximation, with most current coupled climate models adopting a form of the viscous-plastic or elastic-viscous-plastic rheology (Hibler 1979, Hunke and others (1997); see Tandon and others 2018, Hutter and others 2022, Zhang and others 2024 for tabulation of the sea ice schemes in modern Earth system models).
Such models successfully reproduce many aspects of large-scale sea ice dynamics and thermodynamics (Hutter and others 2019, Zhang and others 2011), but are not well suited to the scales of motion observed in MIZs (Herman and others 2022). 
Observations of scale dependence in sea ice deformation using SAR-derived Lagrangian ice trajectories (Marsan and others 2004, Stern and others 2009) have motivated the development of alternative approaches, including the elastic anisotropic plastic (Tsamados and others 2013, Heorton and others 2018) and brittle rheologies (Girard and others 2011, Rampal and others 2016, Boutin and others 2011, Olason and others 2022).
Similar scale dependence has been described in buoy data by, e.g., Rampal and others (2008) and by Itkin and others (2017). 
Scale dependence varies spatially (Hutter and others 2018), hence observations are needed from a wide range of locations and sea ice conditions. 
}

\added{
The lack of consistent, joint observations of FSD and ice motion, over a representative range of sea ice conditions also poses a challenge for sea ice model development and validation.}
\deleted{Reconciling the levels of model complexity needed for large-scale climate simulations, local-scale forecasting, and process-level investigations of sea ice physics remains an active area of discussion (Blockley and others, 2020).}
The inclusion of FSD parameterizations in continuum models (e.g., \cite{hunke1997_ElasticViscous}) is seen as a potential path forward for better simulation of small scale thermodynamic and dynamic sea ice processes and wave-ice interactions \citep{zhang2015_SeaIce,horvat2015_PrognosticModel, roach2018_EmergentSea, bateson2020_ImpactSea}.
Such parameterizations provide links between the statistics of the FSD and other sea ice processes, rather than explicitly representing ice floes.
Alternatively, discrete element method (DEM) models include the effects of the FSD on ice motion either explicitly (disks, polygons, e.g. \cite{gupta2024_EddyInducedDispersion, damsgaard2018_ApplicationDiscrete, manucharyan2022_SubZeroSea}) or implicitly through bonded particles \citep{astromHighResolutionFractureDynamics2024}.
\deleted{It is worth noting that the high computational cost of DEMs has historically limited their applicability to large-scale coupled Earth system models.}
\added{These models have proven especially valuable for process studies, for example examination of coupled interaction between mesoscale ocean turbulence and sea ice melt processes \citep{gupta2022_RegimesSea} and the development of turbulence under sea ice \citep{brenner2023_ScaleDependentAirSea}.}
\deleted{Nonetheless, improved FSD paramterizations from additional observations have advanced the predictability of sea ice and will likely be important inclusions as computational abilities improve.}
Excepting dedicated experiments where multiple observation sources are merged \citep[][e.g.]{uttal2002_SHEBA, hwang2017_WintertosummerTransition, rabe2024_MOSAiCDistributed},  modelers must rely on disparate case studies for the FSD \citep{stern2018_ReconcilingDisparate} and use complementary datasets to validate motion \citep{tsarauHighresolutionRegionalSeaice2024}. 



In this study, we present \added{and analyze novel} observations of \deleted{floe-scale} sea ice \added{properties and} motion \added{at the floe scale} derived from optical satellite imagery. \added{We first provide details of the source data, the floe detection algorithm, and the floe tracking routine. Then, we investigate the statistical properties and seasonal evolution of sea ice floe properties, the FSD, floe-scale motion, and links between motion and floe properties. We compare the floe-scale motion estimates to large-scale ice motion vectors and analyze the drift anomaly distribution. Finally, we consider spatial variation of velocity through analysis of the length scale dependence of sea ice deformation.} 
\replaced{Our focus is on}{We focus on} the Fram Strait and East Greenland \deleted{ice tongue}\added{seasonal ice zone} due to its importance as a connection between the Arctic and the world oceans \citep{smedsrud2017_FramStrait, sumata2022_UnprecedentedDecline}, as it is the main passage of sea ice export from the Arctic Ocean. 
In addition, the presence of sharp gradients in the ocean currents and distinct mesoscale features affecting the ice dynamics \citep{manley1987_EffectsSub, kozlov2021_EddiesMarginal, morozov2023_EddiesArctic, watkins2023_EvidenceAbrupt} make it an ideal testbed for measuring complex, small-scale sea ice dynamics.
Our observations aim to fill a major observational data gap for sea ice dynamics in the spring-to-summer MIZ.
\deleted{Floe shapes detected by the algorithm enable analysis of the sampling variability and seasonality of the FSD.
The observations are presented alongside the results of a discrete element sea ice model simulation showing the effects of mesoscale and submesoscale eddy forcing on sea ice rotation and drift variability.
We emphasize the effects of floe size on sea ice motion, including effects of an evolving FSD, proximity of the sea ice edge, and the presence of a rich upper ocean eddy field.}

\section{Data}
\begin{figure}
\centering{\includegraphics[width=89mm]{fig01_study_region.pdf}}
\caption{Map of the study region (green box) and extent of image details for subsequent figures (\deleted{solid orange, solid purple, and dashed pink boxes}\added{orange and purple boxes}). Blue lines show the mean April sea ice extent for each year 2003-20, while pink lines show the September mean.}
\label{fig:study_region}
\end{figure}

\subsection{Satellite imagery}
We acquired daily 250-m resolution true color (bands 1-3-4) and false color (bands 7-2-1) mosaic imagery from the Moderate Resolution Imaging Spectroradiometer (MODIS) instruments via the NASA Worldview Snapshots utility. 
Two images are used each day, corresponding to the daytime overpasses of the Aqua and Terra satellites. Overpass times for the two satellites are generally between 20 and 90 minutes apart.
Data are provided on a North Polar Stereographic grid.
As imagery is only labeled with the date, we use the Satellite Overpass Identification Tool \citep{hatcher_simon_2022_6475619} to identify the precise overpass time for each image for use in ice motion calculations.
The study region covers the area northeast of Greenland from the Fram Strait in the north to Scoresby Sound in the south (Fig. \ref{fig:study_region}, green box).
We focused on the spring and summer months (April 1 to September 15) to ensure consistent illumination, and cover years 2003-20. Sea ice in the region has a large seasonal cycle and significant year-to-year variability. Blue and red lines in Fig. \ref{fig:study_region} mark the April and September mean sea ice extent for each year from the National Snow and Ice Data Center (NSIDC) Sea Ice Index \citep{fetterer2017_SeaIce}. 
The end-of-season ice extent exhibits greater year-to-year variability than winter extent.

\subsection{Gridded sea ice data products}
We use the NSIDC Daily Ice Motion Vectors data product of time-average sea ice drift at moderate-to-large scales \citep{tschudi2020_EnhancementSea}.
This product merges ice motion estimates from passive microwave imagery, buoys, and from reanalysis wind data.
The motion vectors are provided on a 25-km grid at daily resolution.
Note, however, that the effective resolution is generally coarser than this, as the source satellite motion resolution \added{ranges from} 37.5\added{ to }75 km \citep{tschudi2020_EnhancementSea}. 
We use sea ice concentration (SIC) from the
{NOAA/NSIDC} Climate Data Record (CDR) of Sea Ice Concentration
\added{to characterize the SIC of the region around each floe.}
This product combines SIC estimates from the NASA Team and NASA Bootstrap algorithms by choosing the maximum concentration from each source \citep{meier2021_NOAANSIDC, meier2022}.
\added{While higher resolution SIC estimates are available during some years and seasons (e.g., \cite{spreen2008}), the primary advantage to using the NOAA/NSIDC CDR is consistent, daily time-resolution data across all years and months of the study period, with equivalent accuracy to other satellite-based SIC estimates compared to in situ observations \citep{kern2019_SatellitePassive}.
Due to its design as the maximum of two algorithms, it tends to have a positive bias.
It is important to note that the uncertainty in passive microwave estimates of SIC increases during the melt season and toward the ice edge.
}
\deleted{
We estimate local sea ice concentration at floe positions via nearest-neighbor interpolation.
We calculate the distance to the ice edge by finding the minimum distance between each floe and the set of pixels with zero concentration in the SIC CDR. Similarly, we estimate the distance to the coast by searching for the minimum distance to the grid cells in the NSIDC data products flagged as land.
}
\added{We compute the regional sea ice extent as the area of grid cells in the SIC CDR with SIC $\geq$ 15\%. 
In addition, we compute the MIZ ice extent, defined as the sea ice extent for grid cells with SIC between 15 and 85\%, and the MIZ ice fraction, defined as the MIZ ice extent divided by the regional sea ice extent.}

\section{Methods}
\subsection{Sea ice segmentation and tracking}
We obtain measurements of floe-scale sea ice motion derived from moderate resolution optical satellite imagery using the Ice Floe Tracker (IFT) algorithm \citep{lopez-acosta2019_IceFloe}. \deleted{The study region covers the area northeast of Greenland from the Fram Strait in the north to Scoresby Sound in the south (Fig. \ref{fig:study_region}, green box).} The dataset consists of segmented satellite imagery, floe shape statistics, and derived floe trajectories. We first describe the source data, algorithm, and post-processing routines, then summarize the contents of the dataset. The full dataset is available \deleted{at}\added{from} the Brown Digital Repository at \texttt{https://doi.org/10.26300/ydtf-e778}.
\subsubsection{Ice Floe Tracker algorithm}
The IFT algorithm identifies and tracks ice floe shapes in an image sequence to reconstruct their motion trajectories.
The algorithm consists of five subroutines: (1) data acquisition, (2) pre-processing, (3) \added{segmentation}, (4) tracking, and (5) post-processing (Fig. \ref{fig:flowchart}). Steps 2-4 are described in detail in \cite{lopez-acosta2019_IceFloe} \added{and} \cite{lopez-acostaSeaIceDrift2021} and are briefly summarized here\added{, while the post-processing routine is introduced here.}

\begin{figure}
\centering{\includegraphics[width=89mm]{fig02_flowchart.pdf}}
\caption{Overview of the Ice Floe Track processing pipeline.}
\label{fig:flowchart}
\end{figure}

\begin{figure*}[!ht]
\centering{\includegraphics[width=179mm]{fig03_algorithm_example.pdf}}
\caption{Top: True color (a), false color (b), and processed (c) images for 2013-04-24 from the \textit{Aqua} satellite. In the processed image, masked regions (including the dilated land mask and the cloud mask) are black. Candidate segments are labeled as either floes (orange) or non-floes (light blue); non-floes are discarded in the final product. \added{Middle: As in a-c for the \textit{Aqua} image for 2014-05-01.}
Bottom: Details of \deleted{\textit{Aqua}} true color images from (d) 2013-04-24 and (e) 2014-05-01 with overlaid IFT segmentation results, with objects labeled as floes marked in orange and discarded objects in light blue. Manually identified floe boundary for the true positive floes are shown in dark blue. \added{Bounds for the images are shown in panels (c) and (f), respectively).}
Panel (i) shows the area of manually labeled well-segmented objects versus IFT results for the 2013-04-24 and 2014-05-01 initially (gold) and after area adjustment (blue). Objects with poor correspondence to manual floe labels are marked in gray. The one-to-one line  is shown in black. Image locations for panels (a) through (f) are shown in Fig. \ref{fig:study_region}.} 
\label{fig:algorithm}
\end{figure*}

\deleted{Image}Pre-processing \added{of MODIS true and false color images} comprises de-noising via anisotropic diffusion filtering, contrast adjustment by adaptive histogram equalization and  unsharp masking, and noise removal by morphological reconstruction. These procedures produce a sharpened grayscale image (Fig. \ref{fig:algorithm}\added{a-c, d-f}). Next, we apply a combination of morphological operations and $k$-means clustering to segment the pre-processed image into ice and water regions. \deleted{Watershed segmentation is applied both to the sharpened image and to the segmented image, then these watershed outputs are compared to determine a final segmentation.} \added{Watershed segmentation is used to separate ice regions into individual floes. In this method, images are analyzed as topographic maps by identifying markers,  using local maxima of the distance to water pixels, and producing closed contours by so-called flooding of the image. The method is applied both to the binary $k$-means result and to the sharpened grayscale image. The intersection of the two watershed boundaries is used to divide the $k$-means result into individual floes.}

Finally, we label the connected components of the segmented image and extract shape properties (e.g., area, perimeter, centroid, \added{average reflectance}) from potential floes. The error in discrete representations of shapes tends to increase as the size decreases. In addition, prior studies have found few floes larger than 30-50 km \citep{stern2018_SeasonalEvolution, stern2018_ReconcilingDisparate}. Hence, only candidate floes with \deleted{area} \added{sizes} of at least 300 pixels and at most 90,000 pixels \added{are retained, equivalent to }\deleted{(}18.75-5,625 km$^2$\deleted{;} \added{in area, corresponding to} length scales \added{of} 4.3-75 km\deleted{)are retained}. 
\added{We estimate local sea ice concentration at floe positions via nearest-neighbor interpolation.
We calculate the distance to the ice edge by finding the minimum distance between each floe and the set of pixels with zero concentration in the SIC CDR.
Similarly, we estimate the distance to the coast by searching for the minimum distance to the grid cells flagged as land in the CDR.}

\begin{figure*}[!ht]
\centering{\includegraphics[width=179mm]{fig04_tracked_floes.pdf}}
\caption{Floe shapes and trajectories for a set of floes tracked in \textit{Aqua} imagery from April 27-30, 2014. \added{To reduce visual clutter, floe trajectories are only shown for floes that appear in at least three of the four images}. Satellite overpass time is marked in the title. Image location is shown in Fig. \ref{fig:study_region}. Markers show the position\added{s} of the floe centroid \deleted{at}\added{between April 27 and} the image time\deleted{, while hollow circles show the position of the centroid in the previous images}.}
\label{fig:tracked_floes}
\end{figure*}

The floe tracking component links floe shapes across image pairs\deleted{(Fig. \ref{fig:tracked_floes})}. For each floe, candidate matches in subsequent images are selected by first filtering by travel distance and floe area. Potential floe matches are rotated until the area difference is minimized. If the area difference is sufficiently small after rotation, $\psi - s$ curves \citep{kwok1990_IcemotionTracking} are calculated for each floe. These curves summarize the tangent angle, $\psi$, and arc length, $s$, of each floe. If the $\psi - s$ correlation is higher than 0.7, the floes are linked. Whether a floe can be tracked depends mainly on cloud cover, whether the floe remains intact, and on the algorithm accuracy in identifying the floe shape. \added{As an illustration, a subset of detected floes is shown in Fig. \ref{fig:tracked_floes}, highlighting those which were successfully tracked for at least three of the four depicted days. 
The tracker allows time gaps of at most three days, hence a floe missed in one observation may still be successfully tracked in subsequent images.} \added{Given that up to two position estimates are available for each day (corresponding to the \textit{Aqua} and \textit{Terra} overpass times),}
daily rotation rates are calculated using the mean daily rotation across the two satellites when both are available, otherwise we use rotation from a single satellite. Rotation rates are then normalized to radians per day.

\deleted{Floe positions are regridded to local noon, which is typically less than 2 hours away from the two satellite overpass times.
Sea ice motion in the IFT is derived from matched pairs of ice floes across multiple days.
Thus, as is the case for other remote sensing-based ice drift estimates, we can only measure net displacement between acquisition times, rather than measuring the instantaneous floe velocity. 
The mean floe speed during a day is bounded below by the net displacement per day.
}
\deleted{Ice motion vectors $u_{IFT} = (u_{IFT}, v_{IFT})$ are estimated from the floe centroid time series with reference to the polar stereographic projection.}
\deleted{Up to two position estimates are available for each day corresponding to the \textit{Aqua} and \textit{Terra} overpass times.
We regrid the centroid positions to a fixed-time grid using bilinear interpolation, and estimate the motion using forward differences.}

\subsubsection{Post-processing}
The IFT image processing returns significantly more candidate ice floe shapes than it is able to track. These candidate ice floe shapes contain both untracked floes and false positives, such as ice filaments and clouds. In \cite{lopez-acosta2019_IceFloe}, the floe tracking served as a filter to remove objects such as bright clouds that had been mislabeled as sea ice. In order to leverage shapes from untracked floes for FSD analysis, we developed a post-processing routine to identify which of the untracked objects are likely to be ice floes.
The light blue regions in Fig. \ref{fig:algorithm}d and e are examples of rejected shapes.
We implemented a logistic regression classifier to assign each object a probability of being an ice floe. Training and testing the classifier required identifying a set of true and false positives. True positives were selected from the set of tracked floes such that (a) trajectory mean drift speeds are less than 1 m/s and (b) trajectory maximum drift speed is larger than 0.05 m/s, i.e., moving at least two pixels per day. False positives were identified using circularity and solidity  thresholds. 
Circularity is defined as 
\begin{linenomath*}
\begin{equation}
    C = \frac{4\pi A}{P^2}
\end{equation} \label{eqn:circ}
\end{linenomath*}
where $A$ is the floe area and $P$ is the floe perimeter, such that a perfect circle has $C=1$, while solidity $S$ is the ratio of the area of the shape to the area of the smallest enclosing convex polygon. Objects with $S \leq 0.4$ or $C \leq 0.2$ were labeled as false positives. 

A final source of false positives comes from estimating the local SIC.
We applied nearest neighbor interpolation to the NSIDC SIC data to find the sea ice concentration in the nearest grid cell to the IFT floe positions. Objects with SIC=0 were classified as false positives.
Prior to model fitting, we stratified the sample by year and month, randomly selecting up to 1,000 true and false positives from each subgroup. The randomly sampled data were split into training and testing portions, then trained using 10-fold cross validation. Model fitting was carried out using the Python Scikit-Learn machine learning library \citep{Pedregosa2011}. The algorithm performs well, with a precision score of 0.92 and recall score of 0.9.

Comparison of the IFT floe shapes with manually annotated imagery shows that the methods used to separate neighboring floes produces a systematic underestimate of floe area (Fig. \ref{fig:algorithm}d-f). The detected floe boundary is an approximately constant number of pixels inside the true floe boundary. As a result, the relative error in area is greater for small floes than for large floes. For a floe with detected segment area $A_0$, we implement a simple bias correction factor by adding a fixed offset $\delta L$ to the floe length scale $L_0 = \sqrt{A_0}$, such that the corrected area $A$ is given by
\begin{linenomath*}
\begin{equation}
    A(L_0) = (L_0 + \delta L)^2 \label{eqn:area_adj}
\end{equation}
\end{linenomath*}
and the corrected length scale is $L = L_0 + \delta L$.

To find an optimal value for $\delta L$, we manually labeled $n=1,071$ true floe boundaries for IFT candidate floes across five images. Pairs of true floes and candidate floes were divided into test and training sets, with 2/3 going into the training set. We then identified quality candidate objects using the area fraction of the overlap. Setting thresholds of overlap area to at least 90\% of the candidate floe area and 40\% of the ground truth floe area removed under-segmented objects well, yielding $n=440$ well-segmented floe pairs. Minimizing the mean absolute error in area using the well-segmented pairs yields an optimal correction factor of $\delta L = 8$ pixels, with a mean absolute error of 10.6 km$^2$ (median 4.9 km$^2$) for the well-segmented pairs. Finally, using all test-set floe pairs marked as clean by the logistic regression, absolute error has a long tail with a few large outliers, resulting in mean absolute error of 79.6 km$^2$ (median 12.0 km$^2$).
% In the following, we describe the displacement over time using velocity units (m/s), keep in mind however that in all cases we are discussing the displacement per day.

\subsubsection{IFT data summary}
Detected sea ice floes cover the region of the Greenland Sea where sea ice is typically found \added{(Fig. \ref{fig:data_avail}a-b)}.
The lowest detection frequency is near the edge of the seasonal maximum ice extent, and the highest floe concentrations are seen nearer to the center of the ice tongue\deleted{(Fig. \ref{fig:data_avail}a-b)}.
The highest concentration of detected floes is in the southern portion of the study area. Trajectories reflect the southward, nonlinear motion of sea ice (Fig. \ref{fig:data_avail}c).
The raw segmented imagery contain approximately $10^5$ objects per year, with about half as many passing the quality control algorithm (blue and orange lines in Fig. \ref{fig:data_avail}d).
The number of detected floes passing quality control shows a distinct climatology (Fig. \ref{fig:data_avail}e) with many more floes available early in the year. The number peaks in April and has a secondary peak in June, related to landfast ice breakup. By July, ice extent is much lower and few images have more than 100 floes.
Detecting rotation and drift requires pairs of images and successful floe matching, resulting in the number of rotation estimates being on the order of $10^3$ per year; since an individual trajectory can have multiple rotation measurements, the number of individual tracked objects is slightly reduced (pink and green lines in Fig. \ref{fig:data_avail}d). Most trajectories are between two and ten days long, and the number of trajectories has strong inter-annual variability (Fig. \ref{fig:data_avail}f).


\begin{figure*}
\centering{\includegraphics[width=179mm]{fig05_data_availability.pdf}}
\caption{Top: spatial histograms of (a) all available sea ice floe shapes and (b) tracked floes. Floe centroids were binned into 25 km by 25 km grid cells for enumeration. (c) Ice floe trajectories for floes tracked for a period of at least 7 days.
Bottom: (d) Number of floes per year in the raw IFT output (blue) and after postprocessing (orange). The number of tracked floes is shown in green and the number of rotation rate estimates is in pink. (e). Distribution of the number of floes per image as a function of day of year. Counts were smoothed using a 15-day centered median. (f) Histograms of trajectory lengths in days for each year.}
\label{fig:data_avail}
\end{figure*}

\subsection{Floe size distribution}
The sea ice FSD has been shown to be approximately linear in log-log space for at least some range of floe sizes. 
As such, it has frequently been summarized through fitting a power law of the form in Equation \ref{powerlaw} \citep{rothrock1984_MeasuringSea, stern2018_ReconcilingDisparate}.
A common approach to estimating power law fits of this form from empirical data is to apply a log-transform to the data and fit a line using least squares.
Due to the effects of sparse but high-impact fluctuations in the tail of the distribution, this method can be inaccurate; furthermore, the least squares method does not test whether a power law is appropriate \citep{clauset2009_PowerLawDistributions, bauke2007_ParameterEstimation, goldstein2004_ProblemsFitting}.
Other distributions, such as a truncated power law or lognormal distribution, may fit the data better, while still explaining the approximately linear scaling seen in the log-log plots \citep[e.g.,][]{toyota2006_CharacteristicsSea, herman2010_SeaiceFloesize, burroughs2001_UppertruncatedPower,montiel2022_TheoreticalFrameworka, lu2008_AerialObservations}.
In this study, we follow the approach of \cite{clauset2009_PowerLawDistributions}, as implemented in the Python \texttt{powerlaw} package \citep{alstott2014_PowerlawPython}.
We focus on the power law, truncated power law, and lognormal distributions due to their frequent use in the literature. Multiple distributions are referred to as truncated power laws across different studies. Here, we use the form
\begin{linenomath*}
\begin{equation}
    p(x) = cx^{-\alpha}e^{-\lambda x} \label{truncpowerlaw}
\end{equation}
\end{linenomath*}
where $p(x)$ is the probability of a floe having area $x$, $c$ is a normalizing constant, $\alpha$ is the power law slope parameter, and $\lambda$ is a parameter controlling the rate of cut-off as $x$ grows large. The lognormal distribution is defined using parameters $\mu$ and $\sigma$ and has the probability density function (PDF)
\begin{linenomath*}
    \begin{equation}
        p(x) = \frac{1}{x\sigma\sqrt{2\pi}}\exp\left(-\frac{(\ln x - \mu)^2}{2\sigma^2}\right).
    \end{equation}
\end{linenomath*}
Best-fit parameters for each distribution are identified using maximum likelihood.

\begin{figure*}
\centering{\includegraphics[width=179mm]{fig06_fsd_setup_figures.pdf}}
\caption{Ensemble mean (solid line), standard deviation (dark shading), and min-max range (light shading) of $D_{KS}$ for the best-fit parametric distribution as a function of $x_{min}$ for (a) power law, (b) truncated power law, and (c) lognormal distributions.  Vertical dashed lines mark $x_{min}=$41 km$^2$. (d) Relative error in power law $\alpha$ estimate by subsample size relative to images with greater than 300 floes. Shading as in (a-c).}
\label{fig05_fsd_setup}
\end{figure*}

Since the fitted distribution\added{s} depend\deleted{s} on the minimum value of $x$, we apply a constant $x_{min}$ for the full analysis.
The optimal $x_{min}$ is determined by a simple numerical experiment (Fig. \ref{fig05_fsd_setup}).
Using 300 randomly selected images, for each value of $x_{min}$ from 1 to 75 km we find the best-fit parameters for each parameteric distribution.
The Kolmogorov-Smirnov distance $D_{KS}$, defined as the supremum of the distance between the data and the best fit parametric distribution, is used to measure the goodness of fit.
The optimal $x_{min} = 41$ km$^2$ is found by minimizing the ensemble mean $D_{KS}$ (dashed line in Fig. \ref{fig05_fsd_setup}a-c).
The model fit improves rapidly as $x_{min}$ increases to 41 km$^2$, then degrades past that value.
\added{Variation across ensemble members is small compared to the variation with $x_{min}$ (Fig. \ref{fig05_fsd_setup}a-c).}

It is important to note that the IFT dataset provides a sample of floes in each image, not a complete enumeration.
The identifiable floes are constrained by the characteristics of MODIS data, including its 250 m spatial resolution and frequent cloud cover that obscures many floes in the visible imagery.
Even though the post-processing performs well, some false positives remain in the floe library, adding further uncertainty to the estimated FSD.
The MODIS images and the bounds of the study domain are sufficiently large such that the maximum observable floe size is not constrained by image dimensions, meaning finite size effects \citep{stern2018_ReconcilingDisparate} have a minimal impact on the tail of the recovered distribution.
However, our ability to estimate the FSD in a scene based on a sample of floes depends on the sample size. In the absence of large sets of manually annotated imagery, we simulate the effect of subsampling by computing the relative error in power law $\alpha$ estimates  \added{(Fig. \ref{fig05_fsd_setup}d)}.
We randomly selected 100 images with at least 300 floes in each, then repeatedly sampled $n$ floes from each image without replacement and recalculated $\alpha$. We chose the minimum threshold of floes per image for FSD analysis based on a pragmatic balance between the number of floes available in a typical image by day of year (Fig. \ref{fig:data_avail}e) and the estimated sampling uncertainty (Fig. \ref{fig05_fsd_setup}d).
With $n=100$, we have $\alpha$ errors within $\pm 5\%$ in 80\% of cases, while still retaining some late-summer data.

In addition to the image-by-image calculation, we also calculate FSD fits for the data binned by day of year. 
This estimates a central tendency of the distribution, and enables us to estimate FSD later in the year when the per-image sample size is too low.
\added{In addition, it allows us to examine variation within subsets of the data.}

We weigh the suitability of the three candidate parametric distributions for describing the observed FSD.
The distance between the fitted and empirical distributions was assessed with the K-S statistic $D_{KS}$. 
P-values are determined by repeatedly drawing samples from the fitted distribution, calculating the fit to the sampled values, and calculating $D_{KS}$.
The p-value is then determined by the percentage of the K-S statistics from the resampling procedure that are larger than the $D_{KS}$ for the original data. 

\subsection{Sea ice drift}
\added{Sea ice motion in the IFT is derived from matched pairs of ice floes across multiple days.
Thus, as is the case for other remote sensing-based ice drift estimates, we can only measure net displacement between acquisition times, rather than measuring the instantaneous floe velocity. The mean floe speed during a day is bounded below by the net displacement per day.
Ice motion vectors 
\begin{linenomath*}
    \begin{equation}
        \mathbf{u}_{IFT} = (u_{IFT}, v_{IFT})       
    \end{equation}
\end{linenomath*}
\added{where $u_{IFT}$ is the zonal and $v_{IFT}$ the meridional component of motion} are estimated from the floe centroid time series with reference to the polar stereographic projection.
Up to two position estimates are available for each day corresponding to the \textit{Aqua} and \textit{Terra} overpass times.
The satellite overpass times are typically less than 2 hours away from 12 UTC.
We regrid the centroid positions to a fixed-time grid using bilinear interpolation, and estimate the motion using forward differences.}

\added{For a portion of the analysis, we consider the velocity variability after removing the effects of the mean flow. With reference to an area and time mean $\langle \mathbf{u} \rangle_{\chi, \tau}$, where $\chi$ is the spatial scale and $\tau$ the time scale,} 
the sea ice drift \deleted{ vector $\mathbf{u}$ can be decomposed into an along-track (longitudinal) component $\mathbf{u}_L$ and a cross-track (transverse) component $\mathbf{u}_T$ as: $\mathbf{u} = \mathbf{u}_L + \mathbf{u}_T.$}
\added{anomaly vector}
\added{
\begin{linenomath*}
\begin{equation}
\mathbf{u'} = \mathbf{u} - \langle \mathbf{u} \rangle_{\chi, \tau} \label{eqn:drift_anom}
\end{equation}
\end{linenomath*}}
\noindent
\added{can be expressed as an along-track (longitudinal) perturbation $\mathbf{u'}_L$ and a cross-track (transverse) perturbation $\mathbf{u'}_T$ such that}
\added{
\begin{linenomath*}
\begin{equation}
\mathbf{u'} = \mathbf{u'}_L + \mathbf{u'}_T \label{eqn:drift_comp}
\end{equation}
\end{linenomath*}
}
\added{with $u'_L$ parallel to $\langle \mathbf{u} \rangle_{\chi, \tau}$ and $u'_T$ perpendicular to $\langle \mathbf{u} \rangle_{\chi, \tau} $\citep{lacasce2008_LagrangianStatistics, gabrielski2015_AnomalousDispersion}}.
The transverse drift component is often termed the \emph{perturbation velocity} in the literature. % REFERENCE
\deleted{Similar to Gabrielski and others 2015,}


We use centered-time averages of the NSIDC Ice Motion Vectors product to represent the mean flow.
\added{We use the NSIDC product rather than the IFT vectors for the spatiotemporal average because the IFT vector locations are not homogeneous in time or space.}
\deleted{however here, we are also examining the anomalies relative to the mean flow.}
\deleted{
Note that}
\added{The} properties of \deleted{$\mathbf{u}_T$} \added{$\mathbf{u}'$} \deleted{likely} depend on the choice of the mean flow 
\deleted{$\langle \mathbf{u} \rangle_\tau$, where $\tau$ is the averaging time scale (e.g., }
\citep{rampal2009_turbulentFluctuations}).\deleted{; here, 
we consider the effects of different time-average windows $\tau \in (5, 10, $ and $15)$ where $\tau$ \deleted{is time in days}. }
\added{We treat the NSIDC grid size as an implicit spatial scale and consider time windows of $\tau=5, 10, $ and $15$ days. 
The symmetry of the resulting drift anomaly distribution is an indicator of whether the mean flow has been removed correctly.}

% We use the NSIDC Ice Motion Vectors rather than the IFT vectors for the spatiotemporal average because the IFT vector locations are not homogeneous in time or space.

\subsection{Strain rate estimation}
Area-averaged strain rates are calculated using Green's theorem from sets of point estimates of sea ice velocities along a polygon boundary (e.g. \cite{kwok2003_SubdailySea, Hutchings2012, rampal2019_MultifractalScaling, Dierking2020}).
Given a set of $N$ ice floe positions $\mathbf{x}$ and motion estimates $\mathbf{u}(\mathbf{x})$, we iterate over all possible triangles formed by subsets of floe positions and retain all triangles with minimum interior angle of 20$^\circ$ or greater.
This procedure yields millions of polygons; since the number of available polygons is bounded above by a combinatorial function, the number of polygons varies widely per year. 
Using all possible polygons, as was done in \cite{itkin2017_ThinIce}, results in considerable overlap between polygons in some regions. Due to the large study area, we find that the degree of polygon overlap varies considerably from small to large scales, which may lead to \added{prescribing} uneven spatial weight\deleted{ing} across scales. Small polygons are generally non-overlapping, and large polygons nearly all overlap. To minimize spatial biases from overlapping polygons, we downsample iteratively to find sets of non-overlapping polygons within logarithmically-spaced length scale bins. In each image, we randomly choose a starting polygon, then iterate through the list of remaining polygons, keeping a polygon if the intersection with previously chosen polygons has zero area (e.g., Fig. \ref{fig6_polygons}).
For consideration of spatial scaling, we further downsample by drawing random samples within each length scale bin so that equal numbers of samples within each bin are used.

\begin{figure*}
\centering{\includegraphics[width=178mm]{fig07_polygon_example.pdf}}
\caption{Detected floe shapes in the Terra image for April 24, 2013 (orange). Triangles formed using selected ice floes as vertices (outlined in black) are used to calculate area-averaged strain rates. Examples of triangles at length scales of (a) $17.6$ km, (b) $28$ km, and (c), $48.6$ km length scales are shown, with the length scale defined as the square root of the triangle area.}
\label{fig6_polygons}
\end{figure*}

The Green's theorem approach for strain rate calculation uses the line integral around a region to calculate the area average of the gradient of a quantity.
Given drift vectors $\mathbf{u} = (u, v)$ measured at $N$ points $\mathbf{x} = (x, y)$ forming the vertices of a polygon, the mean velocity gradients within the polygon are calculated as
\begin{linenomath*}
\begin{align}
\overline{u}_x & = \frac 1A \oint u \, dy \nonumber\\ 
& \approx  \frac{1}{2A} \sum_{i=1}^N \left[(u_{i+1} + u_i)(y_{i+1} - y_i)\right],
\end{align}
\end{linenomath*}
where in the summation notation used here, the calculation wraps clockwise around the polygon to the origin. Hence, for the triangles used here, $x_4 = x_1$. Other ice drift gradients are calculated in a similar fashion (see e.g. \cite{lindsay2003_RADARSATGeophysical, Hutchings2018, bouchat2020_ReassessingQuality} for details.) 

From these spatial gradients we calculate the invariants of the strain rate tensor $\dot\varepsilon$ \citep{lepparanta2011_DriftSea}:
\begin{linenomath*}
\begin{align}
    \dot\varepsilon_{I} &= \overline{u}_x + \overline{v}_y \\\dot\varepsilon_{II} &= \sqrt{(\overline{u}_x - \overline{v}_y)^2 + (\overline{u}_y + \overline{v}_x)^2} 
\end{align}
\end{linenomath*}
from which we compute the total deformation
\begin{linenomath*}
\begin{align}
   D = |\dot\varepsilon| &= \sqrt{\dot\varepsilon_{I}^2 + \frac{1}{2}\dot\varepsilon_{II}^2}.
\end{align}
\end{linenomath*}
%%%%% Add the uncertainty estimates for the known
Uncertainty in the strain rate calculations depends on the uncertainty in position $\sigma_x$ and area $\sigma_A$, and on the velocities \citep{Dierking2020}.
Triangle area uncertainty depends on the side lengths, so for a triangle with sides $a, b, c$, the uncertainty in area $\delta_A$ is given by
\begin{linenomath*}
\begin{equation}
\delta_A^2 = \frac{\delta^2}{4}(a^2 + b^2 + c^2)
\end{equation}
\end{linenomath*} % See Fig. 2 in Dierking et al. 2020
and following \cite{Hutchings2012}, relative uncertainty in a strain rate component $\epsilon$
is given by
\begin{linenomath*}
\begin{equation}
    \frac{\delta_\epsilon}{\epsilon} = 2\left(4 \frac{\delta_x^2}{A} + 2 \frac{\delta_x^2}{U^2T^2} + \frac{\delta_T^2}{T^2} + \frac{\delta_A^2}{A^2}\right)^{1/2}
\end{equation}
\end{linenomath*}
From \cite{lopez-acosta2019_IceFloe}, the position error for ice floes detected by the IFT is comparable to the resolution of the MODIS imagery, so $\delta_x \approx 255$ m, and ice drift error $\approx 0.65$ cm/s. 
We considered the effects of minimum interior angle on the strain rate uncertainty by calculating the uncertainty for a right triangle across a range of length scales. Based on this analysis we conservatively choose a minimum angle of 20. Uncertainty decreases with increasing length scale:  relative area uncertainty is less than 5\% for triangles with length scale $L = \sqrt{A_{\text{triangle}}}$ larger than 10 km. 
Uncertainty in the strain rate decreases with increasing triangle area, drift speed, and with decreasing uncertainty in time and space.
With $\delta_T = 1$ hr, $\delta_x = 255$ m, and length scale $15$ km, relative error $\delta_\epsilon / \epsilon$ is less than 1 if the drift speed is at least 1 cm/s, and less than 0.2 if the drift speed is at least 5 cm/s. For the scale analysis, we discard polygons with relative uncertainty greater than 0.5.
Mean relative error for the analyzed length scales decreases monotonically from 0.16 (bin 1, mean $L$ = 18 km) to 0.02 (bin 5, mean $L$ = 67 km).

\subsection{Length scale parameter estimation}
The length scale dependence of deformation is estimated as
\begin{linenomath*}
\begin{align}
    D \sim L^{-\beta},
\end{align}
\end{linenomath*}
% Note: we may be able to analyze time dependence for at least a few days, need to see the number
% of triangles possible.
where $L$ is the length scale of the observation; here, we use $L = \sqrt{A_{\text{triangle}}}$.
Numerous methods have been used to estimate the strain rate scaling parameter $\beta$.
For example, \cite{Marsan2004} and \cite{Hutchings2011} fit a regression line to the logarithm of the bin-averaged total deformation, while \cite{itkin2017_ThinIce} fits the regression function using the full set of log-transformed points.

We analyze the scale dependence using \deleted{both}\added{first} the standard bin-averaging approach, and \added{then} introducing a new approach in which the parameter $\beta$ is treated as an unknown parameter for the distribution of the scaled total deformation. The new approach uses a stratified sample across all length scale bins, rather than fitting to a small number of bin averages. Formally, we hypothesize that the scaled total deformation 
\begin{linenomath*}
\begin{equation}
    D^*(\beta) = D/L^{-\beta} = DL^\beta
\end{equation}
\end{linenomath*}
is lognormally distributed.
Lognormal distributions arise from the central limit theorem when independent positive random variables are combined multiplicatively rather than additively. Lognormal distributions are observed in the plastic deformation of numerous materials \citep{tang2020_LognormalDistributiona, chen2021_WhyLocal} and have been discussed previously in the context of sea ice deformation \citep{Marsan2004}.
Under this assumption, the logarithm of the scaled total deformation is normally distributed 
\begin{linenomath*}
\begin{equation}
    \log D^* \sim \mathcal N(\mu, \sigma^2) \label{eqn:scaled_ln}
\end{equation}
\end{linenomath*}
with parameters $\mu$ and $\sigma^2$ dependent on the parameter $\beta$ and data $D_i^*$ such that
\begin{linenomath*}
\begin{align}
\mu(\beta) &= \frac 1N \sum_{i=1}^N \log D_i^*(\beta) \label{eqn:mu}\\
\sigma^2(\beta) &= \frac 1N \sum_{i=1}^N \left[\log D_i^*(\beta) - \mu(\beta)\right]^2 \label{eqn:sig}
\end{align}
\end{linenomath*}
where $N$ is the total number of triangles and $D_i^*$ is the scaled total deformation of the $i$th triangle. While these are, of course, the standard maximum likelihood estimates for the parameters of a normal distribution, at this stage the $\beta$ parameter needed to calculate $D^*$ is unknown. Our approach is to numerically find the maximum of the log likelihood function for $\mathcal N(\mu, \sigma)$.
The log likelihood function given $N$ observations of deformation $D_i^*$ and length scale $L_i$ is 
\begin{linenomath*}
\begin{align}
    \log \mathcal{L}(\beta | D, L)  = &  -\frac N2 \log (2\pi \sigma^2(\beta))   \nonumber \\ 
    & - \frac{1}{2\sigma^2(\beta)}\sum_{i=1}^N \left[D_i^* - \mu(\beta)\right]^2
\end{align}
\end{linenomath*}

The parameter $\beta$ is then estimated numerically by maximizing $\log \mathcal{L}(\beta | D, L)$ over $\beta \in (0, 1]$ using a stratified random sample of non-overlapping polygons from logarithmically-spaced length-scale bins. We calculate a bootstrap confidence interval for $\beta$ using the quantile method with 1,000 bootstrap replicates.  
The performance of the lognormal log-likelihood approach for estimating the total deformation scaling parameter is evaluated by comparison with the empirical distributions within length scale bins and comparison with the regression-based methods from the literature.

\section{Results}
\subsection{Seasonal evolution of floe properties}
\added{Measured ice floe properties show distinct variation across the spring-to-summer period (Fig. \ref{fig:floe_properties}).
We apply two methods for grouping the floe measurements: dividing floes into pack ice and MIZ based on SIC, and dividing floes into near-coast and near-ice-edge categories.
In the pack ice, median floe length scales decrease from 9 to 8 km from April to July, while the size of the largest floes decreases from 18 to 14 km (90th percentile; top of the light shaded area in Fig. \ref{fig:floe_properties}a). 
Changes in length scale within the MIZ are more muted, with median length scales close to 8 km throughout the spring and summer and large floes beginning to decrease in size after June.
The floe circularity increases from April to June in both the pack ice and the MIZ (Fig. \ref{fig:floe_properties}b).
We note a transition in the seasonal cycle of floe shape occurring near the end of May, after which floes in the pack ice become somewhat less circular, while only minor changes in circularity are seen in the MIZ. 
This occurs at the same time as the secondary maximum in floe count seen in Fig. \ref{fig:data_avail}e.
We hypothesize that the introduction of landfast ice floes results in decreased circularity, as the timing coincides with landfast ice breakup in the Greenland Sea \citep{wadhams1981_IceCover, walsh2022_SeaIce}, and freshly fractured landfast ice tends to have elongated, rectangular shapes.
The timing of the transition also coincides with a darkening of ice floes (Fig. \ref{fig:floe_properties}c), an indicator of ongoing melt and ponding of the sea ice surface.
The average floe reflectance in the MODIS true color imagery is related to but not identical with the broadband albedo of the sea ice; we note that the timing of the change in reflectance is similar to the observed albedo seasonality (e.g., \cite{light2022_ArcticSea}).
For the MIZ, the reflectance increases from April to June, becoming similar to the pack ice reflectance at the start of June. 
Keeping in mind that SIC is used to divide the observations into categories, we note a slight increase in the width of the SIC distribution in the pack ice, with a strong decrease in the median SIC within the MIZ (Fig. \ref{fig:floe_properties}d).}
Noting that the sea ice extent decreases slightly from May to June while the MIZ fraction increases sharply (Fig. \ref{fig:sea_ice_extent}), we suggest that the increase in floe reflectance occurs due to a change in classification of floes from pack ice to MIZ as sea ice concentration decreases.

\begin{figure*}[!ht]
\centering{\includegraphics[width=179mm]{fig08_floe_properties_seasonality.pdf}}
\caption{\added{Statistical summary of sea ice floe properties by day of year subdivided by SIC (top row) and by proximity to boundaries (bottom row). Summaries are only shown for days of year with at least 500 detected floes (summed across all years). In the top row, statistics for the pack ice (SIC $ >$ 85\%) are shown as green lines, with the thick line showing the median, the thin lines the 25th and 75th percentiles, and the dashed lines showing the 10th and 90th percentiles. Statistics for the MIZ (SIC between 15 and 85\%) are shown with purple shading. The thick line shows the median, medium hues show the 25th to 75th percentiles, and light hues show the 10th through 90th percentiles. In the bottom row, floe measurements are divided into ``near coast'' (orange lines, distance to coast less than 100 km and distance to edge greater than 100 km) and ``near ice edge'' (blue shades, distance to edge less than 100 km and distance to coast greater than 100 km). Floe properties organized by column. From left to right: floe size, defined using the square root of the floe area; floe shape, quantified by the circularity (Eq. \ref{eqn:circ}); floe reflectance, defined as the area-average calibrated band 1 (620-670 nm) reflectance from the MODIS true color image; local sea ice concentration, defined as the SIC in the nearest CDR grid cell. Note that in panels a-d, the CDR SIC is the grouping variable so there is no overlap between the pack ice and MIZ SIC distributions.}}
\label{fig:floe_properties}
\end{figure*}

\begin{figure} %[!ht]
\centering{\includegraphics[width=89mm]{fig09_cdr_sea_ice_extent.pdf}}
\caption{\added{(a) Sea ice extent estimated from the NSDIC SIC CDR for the study region (green box in Fig. \ref{fig:study_region}). (b) MIZ fraction, defined as the fraction of the sea ice extent with SIC between 15 and 85\%. In both panels, individual years are depicted with thin line and the 2003-20 average is shown as a thick line.}}
\label{fig:sea_ice_extent}
\end{figure}

\added{
Grouping floes by proximity to the coast and to the ice edge offers an alternative view of MIZ and pack ice evolution. 
Floes near the coast tend to be larger than elsewhere, and we see only small changes in the upper percentiles; near the edge, the length scales are similar to those for the SIC-based definition of the MIZ (Fig. \ref{fig:floe_properties}a).
The prevalence of low-circularity ice floes is higher (Fig. \ref{fig:floe_properties}f), supporting the view of the landfast ice being a source of large floes later in the melt season. 
The length scales of ice floes near the ice edge are distributed similarly to those in the SIC-based MIZ category (Fig. \ref{fig:floe_properties}a,e).
We also note that the circularity increases steadily for floes near the ice edge, in contrast to the SIC-based MIZ grouping (Fig. \ref{fig:floe_properties}b, f).
Floes near the coast are brighter than floes near the ice edge, and both show slight increases in brightness throughout the spring before dropping sharply in June.
Melt onset in the region tends to occur in early May \citep{blissRegionalVariabilityArctic2019}.
The darkening floes in June may be correlated with the emergence of widespread melt ponds on the surface \citep{fengEffectMeltPonds2021, websterSpatiotemporalEvolutionMelt2022}.
The brighter floes near the coast may be due to the known east-west decreasing thickness gradient of sea ice in the Greenland Sea \citep{spreen2020_ArcticSea}.
Finally, we find that SIC decreases within both categories as the season progresses, with initial distributions that are similar to the SIC-based division of pack ice and MIZ, however both distributions overlap such that high concentrations are possible near the ice edge, and late in the season, low concentrations are possible near the coast.
Thus, the statistics of floe shapes, sizes, and brightness are consistent with the expected seasonal cycle of fracture and melt throughout the spring and summer, while showing distinct seasonal patterns within SIC-based and boundary-based categories.
}

\subsection{Floe size distribution}
\deleted{We analyze FSD variability by examining the estimated power-law slope on an image-by-image basis and through binning data across multiple images by day of year (DOY-binning).}
\deleted{Images with suitable numbers of sea ice floes for FSD estimation are found primarily in the early part of the season (Fig. \ref{fig:data_avail}e). 
After mid-July, few images have more than 100 detected floes, so late-summer FSD is only approached with the DOY-binned data.}
\deleted{The decrease in floe numbers over the season is consistent with the known seasonality of the ice cover. 
As the melt season progresses, the ice extent decreases, limiting the area where floes can be found. Furthermore, we expect the size of floes to decrease through the season, as floes fracture and melt. A temporary increase in floe count starting in mid-May and extending until mid-June is observed, consistent with the timing of landfast ice breakup in the Greenland Sea (Wadhams, 1981; Wadhams, 2022).}

\begin{figure*}[!ht]
\centering{\includegraphics[width=179mm]{fig10_likelihood_ratio_results.pdf}}
\caption{Log-likelihood ratios for image-by-image (top row) and DOY-binned (bottom row) FSD parametric distribution fits. In the figure titles, PL=power law, TPL=truncated power law, and LN=lognormal. Log likelihood ratios that are significantly different than zero after accounting for the false discovery rate are colored blue (image-by-image) or green (DOY-binned). Gray dots are not significantly different than zero. Positive values favor the distribution listed first, and negative values favor the distribution listed second, so for example in panel (a) the negative values favor the truncated power law over the standard power law.}
\label{fig:fsd_llr}
\end{figure*}

We next consider the choice of parametric distribution for summarizing the FSD \added{by evaluating log-likelihood ratios \added{(Fig. \ref{fig:fsd_llr})}. This method allows us to determine which candidate distributions are best supported by the observations}.
At a significance level of 5\%, we find that the best fit power-law distribution can be rejected in 23\% of the image-by-image results, which is similar to the results of \cite{hwang2017_WintertosummerTransition}. When binned by day of year, however, the power-law fit can be rejected in 76\% of cases.
\cite{stern2018_SeasonalEvolution} evaluated the performance of the power law for FSD analysis for the Beaufort Sea using the same goodness-of fit test, and found better agreement with the standard power law form than we find here. This may be due to regional differences, or to a difference in methodology--their work uses the mean caliper distance rather than the area to describe variation in floe size. 
The log-likelihood ratio results show that the standard power law is either less likely than the truncated power law or lognormal distribution, or equally likely: only in a handful of images does the standard power law provide a significantly better fit\deleted{ (Fig. \ref{fig:fsd_llr})}. 
Seasonality in the likelihood ratio results suggests a shift in the form of the FSD through the year.
While the sample sizes tend to be lower near the end of the year, consistently large samples from April to June provide confidence that the changes seen during that period are statistically robust.
\deleted{Mokus and Montiel (2022)
%\cite{mokus2022_WavetriggeredBreakup} %
argue that wave-driven breakup leads to a lognormal FSD.}

We find that \added{both the truncated power law and} the lognormal distribution \deleted{is}\added{are} significantly more likely than a standard power law for the DOY-binned data in early summer.
\deleted{, when wave-driven breakup is especially likely to be a contributing factor}
\deleted{It also tends to be more likely than the power law in the image-by-image data. However, }
In most cases, we find that the lognormal and truncated power law distributions are equally likely\deleted{, with a few exceptions}.
After July, the log-likelihood ratio test fails to find any difference in likelihood between the three candidate distributions.
\added{Model experiments by  \cite{mokus2022_WavetriggeredBreakup} suggest that wave-driven breakup leads to a lognormal FSD, while experiments by \cite{zhang2015_SeaIce} and analysis of observations by \cite{herman2010_SeaiceFloesize} support an upper-truncated power law distribution.
We find that the fitted parameters for the lognormal distribution vary widely after mid-June, impeding interpretation, while the parameters for both the standard and truncated power laws are less sensitive to seasonality.}


Examination of the empirical (estimated from the histogram) and fitted PDFs and complementary cumulative distribution functions (CCDFs) for individual dates shows that the truncated power law provides a good fit across most of the data, with larger departures from the fit for very large floes (Fig. \ref{fig:fsd_dates}). 
Gray lines in the figure show the year-to-year variation in the distributions, which can be substantial.
Binning by day of year (black lines) provides an estimate of the central tendency of the distribution: although individual dates may show strong deviations (e.g., Fig. \ref{fig:fsd_dates}d), the binned distributions fall within the range of year-to-year variations.
\added{Given the results of the likelihood ratio tests, the comparisons with the data, and}
\deleted{In order} to facilitate comparison with prior work \added{using slopes to characterize the FSD}, we report results for the truncated power law distribution in the remainder of the study. 


\begin{figure*}[!ht]
\centering{\includegraphics[width=179mm]{fig11_example_dates_fsd.pdf}}
\caption{Empirical (solid) and fitted \added{truncated power law} (dashed) FSD PDFs and CCDFs for randomly selected dates in April (a), May (b), June (c), and July (d). Blue lines show the PDF and CCDF for the selected date, gray lines show the distributions for the same day of year for each year with sufficient data, and black lines show the distribution fitted to binned data for the same day of year. Fitted distributions are only shown for the selected date and for the binned data.}
\label{fig:fsd_dates}
\end{figure*}

\begin{figure}
\centering{\includegraphics[width=178mm]{fig12_fsd_seasonality.pdf}}
\caption{\added{(a)} Fitted FSD slope parameter for the truncated power law distribution by day of year for the image-by-image analysis (blue dots) and for the DOY-binned analysis (red dots). Box-and-whisker plots show the interquartile range and median of the image-by-image data for each month. \added{(b) FSD slope parameter estimated from DOY-binned data separated into pack ice (SIC $> 85\%$, green) and MIZ (SIC between 15 and 85\%, purple). (c) FSD slope parameter estimated from DOY-binned data separated into near-coast (coast distance less than 100 km and edge distance greater than 100 km, blue) and near-edge (edge distance less than 100 km and coast distance greater than 100 km, orange).)}}
\label{fig:fsd_seasonality}
\end{figure}


\added{We find that}
the slope $\alpha$ of the truncated power law distribution shows distinct seasonality amidst strong variability (Fig. \ref{fig:fsd_seasonality}).
For April, May, and June, where there are sufficient images meeting the criteria for the minimum number of floes, we find that the median $\alpha$ increases from 1.9 in April to 2.2 in July. \added{Within individual images, we find high variability in $\alpha$;} the interquartile ranges and interdecile ranges show little variation between months (interquartile range 0.24-0.25, interdecile range 0.46-0.52).  
\added{Images with suitable numbers of sea ice floes for FSD estimation are found primarily in the early part of the season (Fig. \ref{fig:data_avail}e). After mid-July, few images have more than 100 detected floes, so late-summer FSD is only approached with the DOY-binned data.}
Estimates of $\alpha$ based on the DOY-binned data fall within the interquartile range of the image-by-image results.
Binning by day of year enables extension of the analysis into late summer \added{as well as examination of the effects of SIC and spatial variability.}
\added{Across the full study region, }we find that DOY-binned $\alpha$ reaches a maximum near $\alpha=2.5$ between July and August, then decreases thereafter \added {(Fig. \ref{fig:fsd_seasonality}a)}.
For the time range studied here, the minimum \added{in FSD slope} appears to be in April,
\added{coinciding with the early season floe count maximum in Fig. \ref{fig:data_avail}e.}
\added{The observed increase in floe circularity coincident with a preferential loss of large floes is a strong indicator that floe fracture is driving the evolution of the FSD in pack ice during the spring (compare Fig. \ref{fig:floe_properties}b,e).}
\added{Within the pack ice, the seasonal cycle of $\alpha$ mirrors that of the full study region, while within the MIZ, $\alpha$ is higher but does not noticeably vary until July (Fig. \ref{fig:fsd_seasonality}b). 
We see that the slope parameter for the pack ice increases until it exceeds that of the MIZ in late June. In the late season, the MIZ slope parameter is again higher than that of the pack ice.
Near the coast, the FSD slope steepens gradually from April to mid July, while the patterns for near edge and MIZ categories are similar (Fig. \ref{fig:fsd_seasonality}c).
}

The increase in power law slope through the spring-summer melt season has been previously observed in, e.g., the Beaufort Sea \citep{hwang2017_WintertosummerTransition, stern2018_SeasonalEvolution}.
The seasonality we observe here agrees well with a  previous study in the Beaufort and Chukchi seas \citep{buckley2024_SeasonalEvolution}, in which MODIS imagery was analyzed with a different segmentation algorithm and \added{it was} found the running-mean $\alpha$ to range from 1.74 to 2.0 during summer, with minimum in April and maximum in August. 
The timing of the minimum and maximum power law exponent is remarkably similar to the evolution of the FSD \added{for large floes in}\deleted{for} the Beaufort Sea analyzed in \cite{stern2018_SeasonalEvolution}, despite the methodological differences, with minimum $\alpha=1.9$ in April and maximum $\alpha=2.8$ in August.

The pronounced variability in FSD slope is primarily influenced by the strong spatial variability \added{of sea ice and ocean conditions} observed in the Fram Strait.
The position \added{of ice floes} relative to the \deleted{coast line}\added{coastline} (or fast ice edge) and the position relative to the sea ice edge are especially important. 
Landfast ice on the Greenland coast provides a source of large floes as it breaks up through the year. 
Near the ice edge, wave action increases the likelihood of floe fracture, and contact with the relatively warm waters of the North Atlantic enhance local melt rates.
\deleted{Sea ice concentration is higher in the north, where ice advected along the Transpolar Drift Stream enters the Greenland Sea through the Fram Strait.
We examined the dependence of the FSD on SIC and on the distance to the ice boundary (coast and open water). 
While we found some evidence of steeper FSD slopes (higher prevalence of small floes) in low-SIC regions and near the coast, detailed examination of the spatial variability of FSD slope is beyond the scope of this study.}
\added{In summary, we see that the seasonal variation of the FSD includes (a) the transition from predominantly pack ice to predominantly MIZ, (b) spatial heterogeneity with a decaying landfast ice edge, (c) floes becoming smaller and more circular, indicating breakup, and (d) a shift toward an increased role of melt processes in June.
}

\begin{figure} %[!ht]
\centering{\includegraphics[width=89mm]{fig13_all_floes_v_tracked_FSD.pdf}}
\caption{Empirical FSD PDF and CCDF for all floes (\added{both tracked and untracked,} solid black line) and for tracked ice floes only (dashed red line). Thick lines show the distribution across all years, while thin lines show the distribution of individual years.}
\label{fig:fsd_tracked}
\end{figure}

It is worth noting that the analysis thus far has been on the FSD of all detected ice floes, not only the ice floes that were successfully tracked. \added{In Fig. \ref{fig:fsd_tracked}, we compare the FSD for tracked floes (red) with the FSD of all floes (both tracked and untracked, black).}
We see that the set of tracked floes tends to have fewer small floes and a greater number of large floes relative to the set of all detected floes. This is unsurprising, as larger floes are more likely to have better resolved distinguishing features and shapes than do floes near the detection limit.
Differences \deleted{in}\added{between} the distributions \added{of untracked and of tracked floes} in individual years (thin lines) are more extreme than the differences across all years. Analysis of interannual variability will need to account for differences in the detected floe sizes, however for the remainder of this study, we pool data from multiple years together so the results are unlikely to be affected by the differences seen here. 

\subsection{Sea ice motion}
\deleted{We compare the IFT observations to the NSIDC ice motion vectors by binning the observations into 25 km bins (approximately the maximum spatial resolution of the NSIDC product) and by calculating the monthly means across all years (Fig. \ref{fig:mean_drift}). Means are calculated only for times and locations where both IFT and NSIDC data are available. 
Only bins with at least 30 observations are included in the comparison.}
\added{Monthly mean IFT drift vectors show a predominantly southward drift, which decreases in magnitude from April to June (averaged within 25-km bins, Fig. \ref{fig:mean_drift}a-c).
Drift speeds are highest near the ice edge, reaching 20 cm/s in April and decreasing to 12.5 cm/s in June, and lowest near the coast, where the drift speeds drop from 14 cm/s in April to 5 cm/s in June.
}
\deleted{Means are calculated only}\added{Mean drift from the NSIDC Ice Motion Vectors}
for times and locations where both IFT and NSIDC data are available agree in the climatological drift direction but not in the magnitude \added{(Fig. \ref{fig:mean_drift}d-f). 
Drift speeds near the ice edge in the NSIDC product decrease from 11 cm/s in April to 3 cm/s in June, while near the coast the drift speed decreases from 7 cm/s to 2 cm/s in the same period.}

The \added{average} difference between the daily IFT and NSIDC vectors, \deleted{$\mathbf{u}' = \mathbf{u}_{IFT} - \mathbf{u}_{NSIDC}$}
% \added{
% \begin{linenomath*}
% \begin{equation}
%     \mathbf{u}_{diff} = \mathbf{u}_{IFT} - \mathbf{u}_{NSIDC},
% \end{equation}
% \end{linenomath*}
% }
is largely distributed along the climatological mean flow \added{(Fig. \ref{fig:mean_drift}j-l)}.
The mean difference \added{magnitude} shows spatial variability, with the largest vector differences occurring in the southern and eastern edges of the ice tongue. 
Mean differences \added{across the study region} for the months shown are 8.6, 8.9, and 11 cm/s, respectively. 
\added{
Note also that spatial density of retrieved floe motion vectors depends on the sea ice state and thus varies in time (Fig. \ref{fig:mean_drift}g). }
\added{In April and May, most drift vectors are retrieved from the climatological pack ice (SIC $> 85$\% in the 2003-20 mean), shifting to MIZ (SIC between 15 and 85\% in the 2003-20 mean) in June.
Retrieved vector density increases towards the eastern and southern boundaries of the Greenland sea ice tongue.
Following the seasonal cycle of ice retreat (Fig. \ref{fig:sea_ice_extent}), the region with the highest IFT vector density moves north between May and June.}

\begin{figure*}
\centering{\includegraphics[width=179mm]{fig14_mean_drift.pdf}}
\caption{\deleted{Top: Mean daily drift vector for IFT (red) and NSIDC (black). The length of the vector is proportional to its magnitude. The observation count is shaded in blue. Bottom: Black arrows show the vector difference. The magnitude of the difference is shaded in red. Black contours show the 2003-20 average sea ice concentration for each month.}
\added{(a)-(c): Monthly mean drift vector for IFT binned to 25 km resolution at locations with at least 30 resolved vectors. Arrows show the direction and are proportional to the magnitude in length. Drift vector magnitude is shaded. (d)-(f): Monthly mean drift vectors from the NSIDC Sea Ice Motion Vectors product at the same times and locations as the IFT vectors. (g)-(i): IFT vector counts by bin, with the color scale capped at 100 to show detail. (j)-(l): Direction (arrows) and magnitude (shading) of the vector difference between the IFT and NSIDC vectors. Contours show the the isolines for 85\% and 15\% SIC from the NSIDC CDR SIC. To reduce vector density in the drift and drift difference plots, every second vector in the $y$ direction is shown.}}
\label{fig:mean_drift}
\end{figure*}

\deleted{Considering the longitudinal and transverse components of the drift vectors, we see larger variance in the longitudinal direction than in the transverse direction for both IFT and NSIDC. 
In both cases, the variance for IFT is larger than NSIDC (Fig. \ref{fig:drift_dist}). Co-variability of the IFT and NSIDC ice drift vector components was measured by Spearman's rank correlation coefficient. Correlation is highest in April ($\rho=0.68$ for longitudinal, $\rho=0.7$ for transverse) and lowest in July ($\rho=0.13$ for longitudinal, $\rho=0.18$ for transverse).}

\added{Uncertainty in gridded passive microwave motion products arises from numerous sources, including effects of changing surface properties during the melt season as well as sub-grid-scale velocity variability.
Sub-grid-scale velocity variability can come from variation in winds, floe-floe interactions, and small-scale ocean currents.
The mean drift from IFT and NSIDC agree qualitatively in that both show the same direction of drift, and the magnitude of drift has zonal gradient with the strongest drift speeds near the sea ice edge.}
For the NSIDC ice motion product, \cite{tschudi2020_EnhancementSea} noted that the theoretical limit to the precision of the low-resolution (25 km) passive microwave (PM) derived motion vectors is approximately 7 cm/s, while with the higher resolution PM-derived motion available from AMSR-E produces errors closer to 5 cm/s \citep{meier2006_HighresolutionSeaice}. Uncertainty tends to be higher in summer and in the MIZ. \cite{wang2022_IntercomparisonSatellite} compared the NSIDC data to buoy data in the East Greenland region and found errors of 6.62 km/d (7.7 cm/s). Thus the difference between the NSIDC and IFT ice motion estimates \added{shown in Fig. \ref{fig:mean_drift}j-l} is of similar size to prior estimates of the NSIDC uncertainty.
\added{
Furthermore, the NSIDC data product is known to have a negative bias in summer months relative to drifting buoy data \citep{sumata2014_IntercomparisonArctic}.}

\added{
IFT floe trajectories are produced in a Lagrangian framework, similar to how buoy data is acquired. 
Hence, the agreement in the drift direction and mismatch in the drift magnitude may represent a difference in the type of motion detected.
While the NSIDC product aims to provide the average drift within a grid cell, the IFT (and buoys) provide the drift of one piece of ice, which is may differ from its surroundings.
This is especially relevant for the MIZ, where spatial decorrelation length scales are short.
}

\deleted{Uncertainty in gridded motion products arises from numerous sources, including sub-grid-scale ocean turbulence, which can contribute to velocity variability. In the MIZ, the influence of mesoscale eddies is strong. The Rossby radius of deformation in the Greenland Sea region is approximately 5 km (Nurser and Bacon, 2014) and observations of eddies by Kozlov and Adadzhanova (2021) showed mean diameters of 6-12 km depending on ocean depth. The 25 km-minimum spatial resolution of the NSIDC ice motion vectors is thus coarse enough that unresolved eddies within a grid cell may have a significant influence on the uncertainty of grid-mean motion vector.}

\added{We use the drift anomaly relative to the time-average NSIDC motion (Eq. \ref{eqn:drift_anom}, \ref{eqn:drift_comp}) to examine this sub-grid-scale variability. 
}
\added{The drift anomaly distributions for IFT and NSIDC are approximately symmetric about zero (Fig. \ref{fig:drift_dist}a,c).
The tails of the IFT anomalies are broader than for NSIDC. 
For both data products, we see higher variability for the longitudinal case than the transverse case.}
Co-variability of the IFT and NSIDC ice drift vector components \deleted{was} measured by Spearman's rank correlation coefficient \deleted{ Correlation} is highest in April ($\rho=0.68$ for longitudinal, $\rho=0.7$ for transverse) and lowest in July ($\rho=0.13$ for longitudinal, $\rho=0.18$ for transverse). 
Since the distribution of $\mathbf{u}'$ shows little dependence on the length of the time averaging window (Fig. \ref{fig:drift_dist}a, d), we focus on the case with $\tau=5$ days. 

\added{Gaussian distributions of velocity anomalies have been observed for ocean floats \citep{zhang2001_IsopycnalLagrangian}.}
If $\mathbf{u}'$ was Gaussian, then the scaled distribution would be equivalent to a normal distribution with standard deviation 1, shown with a dashed line.
In contrast, \added{analysis of sea ice velocity anomalies in the central Arctic by \cite{rampal2009_turbulentFluctuations} revealed} an exponential distribution PDF. \deleted{shows linear decay in a log-scaled plot (Rampal and others, 2009).} 
\added{
Scaling the IFT drift longitudinal and transverse drift anomaly distributions by the standard deviation, we find that extreme values are more common than one would expect from a Gaussian distribution.
Linear decay of the drift anomaly tails is observed in the pooled data across all months for both the pack ice and the MIZ (Fig. \ref{fig:drift_dist}b, e) and in individual months (not shown).
}
\deleted{We explore the possibility that ocean eddy activity contributes to the variability of $\mathbf{u}'$ using the DEM simulations.}
\deleted{If the sub-grid variability were to arise due to ocean eddy activity, then we would expect that the variability of $\mathbf{u}'$ would be similar to the variability of ice floes in an idealized eddy-rich ocean field. Since the distribution of $\mathbf{u}'$ shows little dependence on the length of the time averaging window (Fig. \ref{fig:drift_dist}a, d), we focus on the case with $\tau=5$ days. The observed $\mathbf{u}'$ distribution scaled by the standard deviation produces a wide-tailed distribution that is very similar to the results of the simulation (Fig. \ref{fig:drift_dist}b, e)).
}
The dotted line in Fig. \ref{fig:drift_dist}b and e shows the best-fit exponential distribution from a random sample of 1000 observations for both drift components (exponential distributions are only defined for positive values; we fit the function to the absolute values and scaled by 0.5 so the PDF normalizes to 1). 
Using a KS-test and bootstrapping on the remaining observations, we find no significant difference between the exponential distribution and the observed drift perturbations. 
\deleted{Thus, the simulation results show that purely eddy-driven ice motion can produce the observed scaled drift anomalies distribution.}


\begin{figure*}
\centering{\includegraphics[width=179mm]{fig15_simpler_velocity_dist_option2.pdf}}
\caption{\added{Left column: Empirical PDFs of longitudinal (a) and transverse (d) components of daily ice motion for IFT (red) and NSIDC (black) }\deleted{, and IFT-NSIDC (blue)} relative to the mean \added{NSIDC} flow vector calculated over a centered window of $\tau = 5, 15, $ and $31$ days (solid, dashed, and dotted lines)\deleted{, shown on left}. Positive values of $u'_L$ and $u'_T$ indicate magnitude of the drift vector projection in the direction of the mean flow and 90$^\circ$ to the right of the mean flow, respectively. 
\added{Center column: longitudinal (b) and transverse (e) IFT drift anomaly distributions after scaling by the standard deviation. Green dots indicate pack ice (SIC $>85$\%) and purple dots indicate MIZ (SIC between 15-85\% SIC).}
Dashed lines show an equivalent normal distribution (mean=0, standard deviation 1) and dotted lines show a fitted exponential distribution. 
Note that the fitted exponential is reflected about $x=0$, as the domain of the exponential PDF is strictly non-negative. 
\deleted{Center column: Observed IFT (blue) and simulated (gray) empirical PDFs of the longitudinal (b) and transverse (e) components of the drift anomalies scaled by the sample standard deviation. }
\added{Right column: Standard deviation of IFT longitudinal and transverse drift anomalies}
\deleted{Longitudinal and transverse components were calculated}relative to the 5-day centered average NSIDC ice motion
\added{as a function of length scale (c) and} \deleted{
(f) Standard deviation of IFT drift anomalies} as a function of edge distance \added{(f)}. \deleted{ As the simulations were performed in a homogeneous eddy field, there is no comparison simulation value here.}
}
\label{fig:drift_dist}
\end{figure*}

\added{The observation of exponential variability in the drift anomaly distribution mirrors the finding of \cite{rampal2009_turbulentFluctuations} for the central Arctic, with notably smaller time scales and spatial scales. In particular, they found spatial scales of 200 km and 2.5 months for summer sea ice.}
\added{
We can expect spatial and temporal scales for the central Arctic to be much higher, due to larger internal ice stresses in high ice concentrations, weaker ocean currents, and high correlations with the wind field in comparison with the Greenland Sea \citep{rampal2009_turbulentFluctuations, thorndikeSeaIceMotion1982}.}
\added{Considering sub-grid-scale component of velocity variability, we would expect the motion of ice floes that approach the same size of the grid cell to have lower error, which we see in the scale dependence of $\sigma_{u'}$ (Fig. \ref{fig:drift_dist}c).
Spatial correlation in sea ice motion also depends on the compactness of the sea ice cover: tightly packed floes are constrained to move together, while in lower ice concentrations floes can follow local variations in oceanic and atmospheric forcing independently. 
Dividing the data into pack ice and MIZ, we find stronger variability in the sub-grid-scale motion in the MIZ in both the longitudinal and transverse directions across length scales and as a function of edge distance (Fig. \ref{fig:drift_dist}c,d).
}


\added{
In the MIZ, mesoscale eddies are a strong contributor to sea ice velocity variability.
The first mode of the Rossby radius of deformation in the Greenland Sea region is approximately 5 km \citep{nurser2014_RossbyRadius}, and observations of eddies by \cite{kozlov2021_EddiesMarginal} showed mean diameters of 6-12 km depending on ocean depth.
The 25 km-minimum spatial resolution of the NSIDC ice motion vectors is thus coarse enough that unresolved eddies within a grid cell may be a significant source of uncertainty for the grid-mean motion vectors.}
\added{Previous work with the IFT has indicated strong coupling with the ocean in the region \citep{lopez-acostaSeaIceDrift2021, watkins2023_EvidenceAbrupt}.}
\deleted{Note that model experiments suggest that a stochastic wind field and interacting sea ice floes can also produce realistic non-Gaussian sea ice motion distributions (Shaddy and others, 2025).}
The influence of ocean eddies on sea ice depends on the link between floe size and eddy size \citep{brenner2023_ScaleDependentAirSea, kim2025_CharacterisationSea} \added{as well as the proximity of the sea ice edge \citep{johannessen1987_IceEdgeEddies}}.
\deleted{In both the observations and the simulations, }
\deleted{We find that $\sigma_{u'}$ decreases moderately as the floe size increases (Fig. \ref{fig:drift_dist}c). }
\deleted{
The standard deviation of the sea ice motion components is much higher in the observations than in the simulation.
One contributor to the difference may be the influence of the fluctuating wind field.
Underestimation of ocean turbulence may also be a contributor.
To that end, it is worth noting that the QG model was tuned to turbulence in the Beaufort Gyre (Manucharyan and others, 2022).
Higher eddy kinetic energy in the Greenland Sea may result in larger sea ice drift variability (Wang and others, 2020)}
\deleted{Spatial variability in sea ice motion depends on the local FSD and on its relationship to the region boundaries. 
The QG model simulates a homogeneous ocean, thus the simulations cannot explore the spatial effects. 
However, from the observations (Fig. \ref{fig:drift_dist}f), we see a strong dependence of $\sigma_{u'}$ on the proximity of the sea ice edge. }
\deleted{This suggests that the role of eddies in sub-grid sea ice variability increases toward the ice edge, which is}
\added{The length and edge distance dependence of the standard deviation of IFT drift anomalies is}
consistent with known factors for eddy generation in the MIZ \citep{kozlov2021_EddiesMarginal, piccolo2024_EnergeticsTransfer}.
Note that model experiments \added{with circular floes} suggest that a stochastic wind field and interacting sea ice floes can also produce realistic non-Gaussian sea ice motion distributions \citep{shaddy2025_WinddrivenCollisions}.
\added{In the next section, we examine the ice floe rotation rates, which provide additional evidence of potential eddy influence on ice motion.}
\added{Future work will weigh the contributions of wind and ocean forcing on the sea ice drift and rotation.}

\begin{figure*}
\centering{\includegraphics[width=150mm]{fig16_rotation_rate_distribution.pdf}}
\caption{\added{Top: Distribution of rotation rates (radians per day) for (a) floes in pack ice and (b) floes in the MIZ. Bottom: Rotation rate variance grouped by month and length scale bin for the pack ice (c) and for the MIZ (d). Mean floe length scale in each bin is indicated in the legend. Only months and length scale bins with at least 100 observations are included.}} 
\label{fig:rotation}
\end{figure*}

\subsection{Rotation rates}
The ability of IFT to resolve sea ice rotation is a key advancement for remote sensing of sea ice motion \citep{lopez-acostaSeaIceDrift2021, manucharyan2022_SpinningIce}.
In the algorithm, rotation rates \added{($\Omega$, rad/day)} are available for a subset of tracked floes (those observed by the same satellite on consecutive days).
\added{As expected from the seasonality of floe counts (Fig. \ref{fig:data_avail}e) and the FSD (Fig. \ref{fig:fsd_tracked}), the number of observed rotations for small floes is larger than for large floes by an order of magnitude, and the number of observations for all floe sizes is largest in Apr-June.}
\added{
Because the MIZ fraction is largest later in the year (Fig. \ref{fig:sea_ice_extent}), we tend to have more observations in pack ice than in the MIZ.}
Nonetheless, the number \added{of rotation observations} is sufficient to \deleted{examine the}\added{identify statistically robust} relationship\added{s} between floe size and rotation \added{and, for a subset of months and length scales, to examine the seasonality.}

\added{We find }
\deleted{In both the simulations and the observations, we see }a narrowing of the rotation rate distribution as floe size increases (Fig. \ref{fig:rotation}a,\added{b), consistent with prior findings for the Beaufort Sea \citep{manucharyan2022_SpinningIce}}. 
The highest rotation rates are seen in the smallest floes 
\added{and in the MIZ. The scale dependence is much stronger in the MIZ than in the pack ice (Fig. \ref{fig:rotation}a,b).}
\added{We see a tendency toward anticyclonic rotation, with slight negative skew in the pack ice and a strong tendency in the MIZ. 
The length-scale dependence seen in Figs. \ref{fig:rotation}a,b persists through the seasonal cycle.
Binning the rotation rates by month and length scale, we see that the variance of $\Omega$ is inversely proportional to the floe length scale in each month (Fig. \ref{fig:rotation}c,d).
Some evidence of seasonality is seen, suggesting a maximum variance in June. 
}

\added{Previous work comparing IFT floe rotation rates to ocean models and observed ocean currents demonstrated that the variance of the rotation rate is correlated with oceanic eddy activity \citep{lopez-acostaSeaIceDrift2021, manucharyan2022_SpinningIce, kim2025_CharacterisationSea}.}
Smaller floes have lower inertia and thus respond more quickly to oceanic forcing \citep{manucharyan2022_SpinningIce, kim2025_CharacterisationSea}.
\added{We infer that the narrow peak in pack ice indicates constraints on the rotation rate from floe interactions, while the wider distribution in the MIZ indicates both fewer floe interactions and higher numbers of ocean eddies.}
\deleted{In addition, t}The relationship between ocean vorticity and floe size depends on the relative size scales of ocean eddies and sea ice floes \citep{kim2025_CharacterisationSea}.
The lower rotation rates of large floes may then reflect the ocean currents varying over smaller scales than the size of the floe \citep{brenner2023_ScaleDependentAirSea}. 
\added{In coupled model experiments, feedbacks between ice and ocean dynamics and thermodynamics depend on the relative sizes of ice floes and ocean eddies \citep{gupta2024_EddyInducedDispersion}.
Smaller floes are more easily dispersed by eddies, resulting in enhanced melt rates.
Future work should examine the relationship between the FSDs and rotation rate distributions obtained through IFT and the rate of sea ice breakup and melt.}
\deleted{The observed rotation rate distributions, while showing a similar scale dependence as the simulations, have narrower peaks and wider tails. 
One reason may be that the floe observations are not limited to low sea ice concentrations. 
We hypothesize that the narrow peak indicates constraints on the rotation rate from floe interactions, while the wider tails arise from the more energetic eddy field in the Greenland Sea compared to the simulations.}

\begin{figure*}
\centering{\includegraphics[width=178mm]{fig17_deformation_scales.pdf}}
\caption{Total deformation rate as a function of polygon length scale ($L = \sqrt{A_{triangle}}$) for months with sufficient data. Blue dots show the stratified random sample. Blue lines show the binned percentile-based statistics (10th, 25th, 50th, 75th, and 90th percentile), while black dots show the bin averages. Whiskers on the black dots show the 95\% confidence interval for the mean. Red triangles and dots show the analytical median and mean using the best-fit scaling parameter.}
\label{fig:strain_rates}
\end{figure*}

\subsection{Deformation length scales}
We next examine the spatial variation in sea ice motion by analyzing deformation.
Total deformation $D$ decreases with increasing polygon length scale (Figure \ref{fig:strain_rates}).
Estimated length scale parameters using the MLE method are $\beta=0.67$ in April, $\beta=0.66$ in May, and $\beta=0.51$ in June (bootstrap 95\% confidence intervals are shown in the figure). Values of $\beta$ estimated using the least squares (LSQ) method are slightly higher in each month (0.7, 0.73, and 0.58, respectively), but the differences between the LSQ and MLE methods are not significant.
The decrease in $\beta$ in June is statistically significant for both MLE and LSQ methods. 
It primarily represents an increase in $D$ at length scales above 25 km, with small changes for the smallest length scale bin (mean $L = 18$ km) and an increase from 0.13 day$^{-1}$ to 0.19 day$^{-1}$ at the largest length scale bin (mean $L = 66$ km). Thus, through the melt season, the total deformation rate increases, while the deformation becomes less localized. 
\added{These changes occur while the FSD steepens and the MIZ fraction increases.}
\deleted{The decrease in the localization of the deformation rates occurs during while the FSD steepens.}
Future work should examine the correspondence between the FSD and deformation properties.

The acquired scaling parameters are much steeper than has been observed in the central Arctic winter pack ice (e.g., $\beta=0.21$ in Hutchings and others (2011, 2012, 2018) and $\beta=0.2$ in \cite{Marsan2004}), and are similar to the values found for small length scales during the N-ICE 2015 expedition by \cite{oikkonen2017_SmallScale} for sea ice north of the Fram Strait.
In particular, the N-ICE2015 data show $\beta$ increasing from $0.52$ to $0.82$ as the time interval decreased from 24 hours to 10 minutes, and found lower deformation rates as the distance from the ice edge increased.
Our data is primarily concentrated further south than the N-ICE2015 study location, which may contribute to the difference in observed scaling parameters as the sea ice properties vary with latitude.

The strong agreement between the lognormal MLE and the LSQ results gives confidence in the approach. The analytically derived lognormal mean and median agree very well with the empirical estimates. 
One advantage of the MLE approach is that a larger fraction of the data can be used directly in the estimation, in contrast to a fit to a small number of bin averages used in the LSQ approach.
Assessment of the lognormality of the data with the MLE $\beta$ parameters via qq-plots showed good agreement between the scaled data and the theoretical percentiles of the log-normal distribution.
We also note that the good performance of the MLE fit for the moderate length scales analyzed here does not imply that sea ice deformation is self-similar across all length scales. \added{As shown by \citep{hutchings2024_SeaIce}, sea ice deformation at small scales (meters to 10 km) does not follow the same scaling laws as deformation at large scales (10 - 100 km). Interaction at the scale of individual floes differs from the interaction of aggregations of floes. In future work, extending the IFT observations to include unresolved ice between identified ice floes will allow deeper analysis of the effects of the local FSD on deformation.}

\section{Conclusions}

\added{IFT fills an important gap in the observational network, providing comprehensive data in a season and location where other methods are unavailable or have high uncertainty.}
\added{Leveraging 18 years of MODIS imagery, we show for the first time that the FSD in the Greenland Sea displays distinct, spatially heterogeneous seasonality.
The FSD steepens from the time of initial ice breakup in mid-April to its maximum in July. 
The timing and magnitude of the seasonal cycle depends on the proximity to the coast and to the sea ice edge, and differs between the pack ice and the MIZ.
Observations of floe shape and color contextualize the seasonal variations in FSD.
Steady increases in floe circularity, coincident with preferential loss of large floes, are indicators of floe fracture driving the steepening of the FSD.
Darkening ice beginning in June suggests the formation of melt ponds and an increase in the role of thermodynamic processes for shaping the FSD.}

The floe-scale observations presented here constitute a new class of \added{ice motion} observations for the spring and summer MIZ. \added{By combining information on floe size with translation and rotation, the IFT enables new insights into scale-dependent ice motion in the complex and dynamic MIZ.}
\deleted{complementing existing in situ and remote sensing techniques by increasing the sample size of floe-scale displacement data in the MIZ relative to available drifting buoy observations, by linking floe shapes with ice dynamics, and by producing floe rotation rates.}
\added{Building on our prior work, we provide observational evidence that the variability of both the sea ice drift anomalies and rotation rates are related to the floe size. As floe size increases, the scale of the drift speed anomaly and rotation rate distributions decreases.}
The ability to resolve rotation is a key advancement for remote sensing of ice\added{, with implications for ocean monitoring.}\deleted{, opening new opportunities for research.}
\deleted{For example, the} \added{Ice floe} rotation rate \deleted{data}\added{distributions} \deleted{can be used to estimate} \added{depend both on floe size and} on the underlying ocean vorticity.
\added{Thus, IFT floe-scale observations have particular value for ocean monitoring, for example as data assimilation inputs for ocean state estimates.}
\deleted{Notably, the IFT observations allow simultaneous analysis of the FSD and drift variability.
These observations fill a gap in the observational network, providing data in a season and location where other observations are lacking or have high uncertainty--specifically, during the melt season.
Analysis of the FSD across 18 years of data from the spring-\added{to}-summer transition demonstrates a seasonally evolving FSD.
We investigated the structure of the FSD and found that an exponentially truncated power law distribution was better supported by the data than a standard power law or a lognormal distribution.
These spatially extensive, long-term observations of the FSD provide a new angle to characterize the seasonal evolution of the MIZ, complementing existing remote-sensing derived metrics based on SIC.}
\deleted{Building on our prior work, we provide observational and model-based evidence that the variability of both the sea ice drift anomalies and rotation rates are related to the floe size. In particular, as floe size increases, the scale of the drift speed anomaly and rotation rate distributions decreases.}
% link this to the ocean dynamics
\deleted{The simulations show that the observed scale dependence can arise from ocean turbulence in the absence of winds and large-scale currents.}
\deleted{Future work should weigh the relative roles of wind and ocean forcing in the development of the drift speed anomaly distribution.}

\added{Floe tracking from optical imagery greatly expands the number of Lagrangian motion observations for the Greenland Sea. The low residence time of sea ice floes in the region makes it impractical to regularly deploy buoys, making observations spatially sparse. Thus, the IFT observations fill a major data gap for measuring and understanding spatial variability and scale dependence in sea ice motion.}
This study marks the first time, to our knowledge, that the deformation length scales parameter in the MIZ has been calculated using floe-scale remote sensing observations.
Notably, the deformation scale dependence is stronger than in the central Arctic and varies seasonally.
Deformation rates increase during the spring season, and the scale dependence of the deformation rates varies by month.

\added{These observations open new avenues for characterizing the interannual and seasonal variation in sea ice properties in MIZs.
For example, the existence of a consistent seasonal cycle in the FSD may be leveraged to improve seasonal forecasting of ice conditions.
Future work should examine year-to-year variation in the timing of the onset of FSD steepening and the transition point between fracture and thermodynamic drivers of FSD evolution.}
\deleted{Future work will examine the}\added{ Joint measurements of floe shape and floe motion allow investigation of the} connection\added{s} between area-average deformation properties and granular flow, in particular the potential of identifying differences in deformation rates for strongly and weakly interacting sea ice fields.
\added{We anticipate these observations having high value for development and validation of parameterizations of sea ice fracture, melt, and FSD-dependent dynamics in models.
}
\deleted{A major motivation of the work presented here is to enable model development, including new parameterizations of FSD effects in sea ice models and representation of mesoscale and submesoscale sea ice-ocean interactions. 
In particular, the data presented here is anticipated to provide new avenues for evaluating FSD-dependent ice dynamics in both continuous and discrete element sea ice models.}

\section{Acknowledgements} 
DMW, JKH, and MMW were supported by NASA The Science of Terra, Aqua, and Suomi-NPP Program (20-TASNPP20-0202, 80NSSC22K062). DMW, EMB, MK, and MMW were supported by Office of Naval Research (ONR) Arctic Program (N00014-20-1-2753, N00014-22-1-2741, and N00014-22-1-2722) and the ONR Multidisciplinary University Research Initiatives Program (N00014-23-1-2014).
MK and MMW were supported by the ONR Young Investigator Program (N00014-24-1-2283).

% MODIS acknowledgement
We acknowledge the use of MODIS True Color and False Color Corrected Reflectance imagery from the \textit{Terra} and \textit{Aqua} satellites acquired via the from the Worldview Snapshots application, part of the Earth Observing System Data and Information System (EOSDIS) URL: {\texttt{https://wvs.earthdata.nasa.gov}}. 

\bibliography{summer-ice-velocity}   % reads igsrefs.bib
\bibliographystyle{igs}  % imposes IGS bibliography style on output

% \appendix
% \section{Appendix}

\end{document}